\documentclass[10pt,a4paper]{article}
\usepackage[margin=2.2cm]{geometry}
\usepackage{fontspec}
\usepackage{kotex}
\setmainfont{Apple SD Gothic Neo}
\setsansfont{Apple SD Gothic Neo}
\usepackage{amsmath,amssymb}
\usepackage{graphicx}
\usepackage{caption}
\usepackage[dvipsnames]{xcolor}
\usepackage[colorlinks=true,linkcolor=MidnightBlue,citecolor=MidnightBlue,
            urlcolor=MidnightBlue]{hyperref}
\usepackage{titlesec}
\titleformat{\section}{\large\bfseries}{\thesection}{0.6em}{}
\captionsetup{font=small,labelfont=bf}
\setlength{\parskip}{0.35em}
\setlength{\parindent}{0pt}
\usepackage{mdframed}
\newmdenv[linewidth=0.4pt,linecolor=gray!60,backgroundcolor=gray!5,
          innertopmargin=6pt,innerbottommargin=6pt,skipabove=8pt,skipbelow=8pt]{claimbox}
\newcommand{\claim}[2]{\begin{claimbox}\textbf{주장.} #1\par\smallskip
  \textcolor{gray!70}{\textbf{반증 조건.} #2}\end{claimbox}}

\title{\bfseries 여섯 번째 문패가 닫혔다\\[2pt]
\large 지도는 있고 원인은 모른다 --- 그리고 영화가 venue 가설을 그 자리에서 반증했다}
\author{월드모델 실험실 · 노트 $842$}
\date{2026-08-08}
\begin{document}\maketitle

\section{한 줄}

마지막 문패($822$ --- 합동 이득 지도의 제$3$ 변수)를 닫았다: $n=11$ 스피어만에서
venue 몫 $r=-0.182$ · t\_b 몫 $+0.027$ · 행수 $+0.118$ --- \textbf{전부 문턱
$0.75$ 미달 $\to$ (없).} 정직한 최종 문장: \emph{``합동이 어디서 이기는지는
안다(지도 · $828$ 확정). 왜 이기는지는 아직 모른다.''} $837$ 재기선 후 다섯
사이클 만에 \textbf{문패 $6$건($808 \cdot 813 \cdot 814 \cdot 818 \cdot 819
\cdot 822$) 전량 완결}이다.

\begin{figure}[t]\centering
\includegraphics[width=0.97\textwidth]{figs/six.pdf}
\caption{왼쪽: venue 몫 대 $\Delta$ --- 영화와 아이돌이 정반대. 오른쪽: 문패
결산.}
\end{figure}

\claim{\textbf{영화가 $841$ 의 venue 가설을 그 자리에서 반증했다.} 신규 측정:
챔피언 $0.4862$ 대 자기 릿지 $0.3712$ $\to$ $\Delta = +0.115$(합동 우세 방향 ·
점추정 --- 규약 $47$ 구간 판정은 별도). \textbf{venue 몫 $98.7\%$ 도메인(영화)
에서 합동이 이기고 $93.2\%$ 도메인(아이돌)에선 진다} --- 내부 정렬이 같은 두
도메인이 이득 부호는 정반대다. 소무리 소속만으로는 부호가 안 갈리고, 남는
대비는 학습 행수($481$ 대 $122$)와 라벨 잡음 --- 다음 확장 도메인이 오면
재상관(미결 · 결정 규칙 등록).}
{내 부 예측(`영화 판정 불능 $60\%$')도 틀렸다 --- 영화 릿지는 venue 축을
그대로 갖고도 합동을 못 따라간다. GitHub 흐름 $5$회째 완주(이슈 \#$10 \to$
PR \#$11$). 판 $\rho$ $0.4710$ 그대로.}

\end{document}
